\section{Finite windows: diagnostics and their limits}\label{sec:windows}

The following numbers reproduce reported high-precision computations.
They are not used in the analytic proofs. We do not supply here
end-to-end enclosures for matrix entries, omitted source tails,
eigenvalue isolation and every derived statistic. Accordingly all
printed numerical values in this section are interpreted as
approximations, not as certified bounds on a continuum infimum.

\begin{figure}[t]\centering
\includegraphics[width=.8\linewidth]{figures/fig3_window_bottoms.pdf}
\caption{Reported smallest eigenvalues of finite CCM compressions as
the Fourier cutoff $N$ increases, at $m=\lambda^2=13,23,43$.
The full carrier has $2N+1$ modes. The apparent finite-cutoff
stabilization is a numerical observation, not a proof of equality
with the continuum bottom.}\label{fig:bottoms}
\end{figure}

\paragraph{Finite bottoms.}
For the full compression $K(m,N)$, let $\lambda_1(m,N)$ be its least
eigenvalue. It is an upper bound for the normalized window infimum
$\mu_\lambda$ (\Cref{sec:setup}). The records associated with
\Cref{fig:bottoms} report stabilization near
$N^*/m\approx4.6,6.3,8$ and values approximately
$3.48\cdot10^{-59},1.8\cdot10^{-112},1.7\cdot10^{-220}$.
This does not prove convergence along a varying-window schedule or
provide a lower enclosure for $\mu_\lambda$.
The last printed value is about $46\%$ below $10^{-219.5}$, so it
does not support a claimed agreement within $20\%$.
No cross-implementation identity or asymptotic decay law is inferred
from this comparison.

\paragraph{Which trial is measured.}
The continuum prolate packet $k_\lambda$ of
\cite[\S7]{CCM-ZST-2025} is distinct from its finite projection.
Write $q_{m,N}$ for the unit coefficient row of the projected packet
in the same Fourier basis as $K(m,N)$, and $\xi_{m,N}$ for the unit
lowest row used in the finite computation. All overlaps below refer
to these two finite rows. For a row $v$ with $F_v(0)\ne0$, define
$\kappa(v)=-F_v''(0)/(2F_v(0))$ using its finite P59 transform.
We use
\[
 \kappa_\Xi=\tfrac12(\log\xi)''(\tfrac12)
 =\tfrac12\left[-8+\tfrac14\psi_{\rm d}'(\tfrac14)
                  +(\zeta'/\zeta)'(\tfrac12)\right]
 \approx0.0231049931,
\]
where $\psi_{\rm d}$ is the digamma function.
A formal continuum perturbation calculation suggests, for the
measured diagonal $N=m$, the conjecture
\begin{equation}\label{eq:trialjet}
 \kappa(q_{m,m})=\kappa_\Xi-\frac1{16\pi m}
                     -\frac{13}{256\pi^2m^2}+O(m^{-3}).
\end{equation}
We do not prove this expansion or the finite-projection error bound
needed to obtain it from a continuum calculation. A ratio to
$\Xi(z)/\Xi(0)$ is interpreted only near the nonzero anchor, not
across an unproved cancellation of all zeros of $\Xi$.

The reported values of
$a_m=m(\kappa_\Xi-\kappa(q_{m,m}))$ are
$0.020307,0.020123,0.020016,0.019957$ at $m=13,23,43,83$.
A least-squares fit $a_\infty+\beta_{\rm fit}/m$ gives approximately
$(0.019891,0.00540)$, compared with the proposed
$(1/(16\pi),13/(256\pi^2))\approx(0.019894,0.005145)$.
The discrepancy in the second coefficient is about $5\%$.
These rounded observations do not establish either coefficient or
the stated remainder. The $m=83$ datum concerns the trial coefficient
only and is not a row of \Cref{tab:jet}.

\begin{table}[t]\centering\small
\begin{tabular}{@{}rrrrrrr@{}}\toprule
$(m,N)$ & $T_{m,N}$ & $\kappa(\xi)-\kappa_\Xi$
 & $\kappa(q)-\kappa_\Xi$ & $\delta_{m,N}$
 & $\delta_{m,N}/T_{m,N}$ & $p_{m,N}$\\ \midrule
(13,13) & $1.234\cdot10^{-2}$ & $+2.791\cdot10^{-3}$ & $-1.562\cdot10^{-3}$ & $+4.353\cdot10^{-3}$ & 0.353 & $3.66\cdot10^{-3}$\\
(23,23) & $1.059\cdot10^{-2}$ & $+3.158\cdot10^{-3}$ & $-0.875\cdot10^{-3}$ & $+4.033\cdot10^{-3}$ & 0.381 & $2.98\cdot10^{-3}$\\
(43,43) & $8.236\cdot10^{-3}$ & $+2.738\cdot10^{-3}$ & $-0.466\cdot10^{-3}$ & $+3.204\cdot10^{-3}$ & 0.389 & $1.86\cdot10^{-3}$\\
(13,120) & $1.383\cdot10^{-3}$ & $-1.567\cdot10^{-3}$ & $-1.562\cdot10^{-3}$ & $-4.6\cdot10^{-6}$ & $-0.003$ & $4.69\cdot10^{-9}$\\
\bottomrule\end{tabular}
\caption{Reported rounded diagnostics, with
$\delta_{m,N}=\kappa(\xi_{m,N})-\kappa(q_{m,N})$,
$p_{m,N}=1-|\langle\xi_{m,N},q_{m,N}\rangle|^2$ and
$T_{m,N}=L^2(4\pi^2)^{-1}\sum_{k>N}k^{-2}$.
Rows with small nonzero $p$ are closely aligned directions, not equal
vectors. Small nonzero $\delta$ is not a vanishing theorem. Entries
are rounded separately.}\label{tab:jet}
\end{table}

\begin{figure}[t]\centering
\includegraphics[width=.6\linewidth]{figures/fig4_trial_jet.pdf}
\caption{Reported finite-trial second-jet coefficient $a_m$ against
$1/m$, compared with the conjectural expansion \eqref{eq:trialjet}.}
\label{fig:jet}\end{figure}

\paragraph{The relative statistic and the positivity boundary.}
On cells where $\lambda_1(m,N)>0$, the reported values approximately
$1.35$ and $1.33$ refer to
\[
 \frac{R_{m,N}(q_{m,N})}{\lambda_1(m,N)},\qquad
 R_{m,N}(q)=\frac{\langle q,K(m,N)q\rangle}{\langle q,q\rangle},
\]
at $(13,120)$ and $(23,160)$. They do not refer to
$\|\xi-q\|_2/\sqrt{\lambda_1}$, and they are not denoted by a
continuum-limit subscript. A ratio inequality cannot be multiplied
by an unknown-sign denominator to infer positivity.
If nonnegativity of the true window bottoms were independently
proved cofinally, CCM's monotonicity would imply nonnegativity on
every fixed smaller window and hence the Weil criterion. Proving
such a premise would be legitimate; assuming it would not be a
proof. The finite observations here establish neither that premise
nor a continuum saturation estimate.
